\chapter{Zadatak 1f --- osjetljivost optimalne vrijednosti}
\label{ch:z1-osjetljivost}

\section{Što se ovdje traži}

Odgovoriti koliko se optimalna vrijednost mijenja ako se ograničenja malo
pomaknu --- i to obrazložiti, ne samo navesti broj.

\begin{zadatak}{Zadatak 1, podzadatak f)}
f) Što možete reći o osjetljivosti optimalne vrijednosti ($p^\star =$ opt.\
vrijednost primarnog problema) na perturbacije u ograničenjima (na perturbacije
koeficijenata $b_1$ i $b_2$)?
\end{zadatak}

\begin{zamka}[Prvo razjasniti: što su $b_1$ i $b_2$]
U tekstu zadatka slovo $b$ se prije podzadatka f) nigdje ne pojavljuje ---
ograničenja su zapisana s koeficijentima $a_{ij}$ i brojevima $-10$ i $3$ na
desnoj strani. Iz konteksta su $b_1$ i $b_2$ upravo te desne strane:
\begin{equation}
b_1 = -10, \qquad b_2 = 3, \qquad\text{dakle}\qquad h = \begin{bmatrix} b_1 \\ b_2 \end{bmatrix}.
\end{equation}
,,Perturbacija koeficijenata $b_1$ i $b_2$`` znači pomicanje tih desnih strana,
tj.\ stezanje ili popuštanje ograničenja. To je tumačenje, ne citat --- na
obrani ga izgovori kao svoje tumačenje.
\end{zamka}

\begin{predznanje}
\begin{itemize}[nosep,leftmargin=*]
  \item Optimalni množitelji $\muopt = (404/33,\; 0)$ --- \eqref{eq:z1b-rjesenje}.
  \item Jaka dualnost vrijedi (Slater) --- \S\ref{sec:z1a-slater}.
  \item Perturbacijski vektor $v$ --- \S\ref{sec:oznake-optimizacije}.
  \item Rezerva neaktivnog ograničenja, $g_2(\xopt) = -49/33$ --- \S\ref{sec:z1b-slucaj2}.
\end{itemize}
\end{predznanje}

\section{Perturbirani problem}
\label{sec:z1f-perturbirani}

\begin{izvor}{Predavanje 4 --- Lagrangeova dualnost, slajd 23/26 (PDF str.~31)}
\slajd[0.86]{p04-perturbirani-problem}

\textbf{Transkripcija.} \emph{Perturbiran optimizacijski problem.}
\[
\begin{aligned}
\text{Primaran:} \quad \min_x \ & f(x) \\
\text{uz ogr.} \quad & h_i(x) = u_i, && i = 1, \dots, m \\
& g_j(x) \leq v_j, && j = 1, \dots, p
\end{aligned}
\qquad
\begin{aligned}
\text{Dualan:} \quad \max_{\lambda,\mu} \ & \dual(\lambda,\mu) - \lambda\trans u - \mu\trans v \\
\text{uz ogr.} \quad & \mu \geq 0
\end{aligned}
\]
$u = [u_1 \dots u_m]\trans$, $v = [v_1 \dots v_p]\trans$ su perturbacijski
parametri. Neka je $\popt(u,v)$ opt.\ vrijednost primarnog problema, kao funkcija
od $u$ i $v$; $\popt(0,0)$ je optimalna vrijednost neperturbiranog primarnog
problema.

\textbf{Ključna zapažanja:}
\begin{itemize}[nosep,leftmargin=*]
  \item Perturbacija ulazi na \emph{desnu} stranu ograničenja. $v_j > 0$
        \emph{popušta} $j$-to ograničenje (dopustivi skup raste), $v_j < 0$ ga
        \emph{steže}.
  \item Dualni problem se mijenja samo za dodatni linearni član $-\mu\trans v$ ---
        i baš zato je veza između $\mu$ i osjetljivosti tako izravna.
\end{itemize}
\end{izvor}

Za naš zadatak, s $g(x) = Gx - h$, perturbirani problem glasi
\begin{equation}
\popt(v) = \min\left\{ \tfrac12 x\trans Q x + c\trans x \ \middle|\ Gx - h \leq v \right\},
\label{eq:z1f-perturbirani}
\end{equation}
a to je isto što i pomicanje desnih strana s $h$ na $h + v$, dakle
$b_1 \to b_1 + v_1$ i $b_2 \to b_2 + v_2$.

\section{Globalna granica}
\label{sec:z1f-globalna}

\begin{izvor}{Predavanje 4 --- Lagrangeova dualnost, slajd 25/26 (PDF str.~33)}
\slajd[0.86]{p04-globalna-osjetljivost}

\textbf{Transkripcija.} Iz $\popt(u,v) \geq \popt(0,0) - \lambda^{\star\mathsf{T}}u - \mu^{\star\mathsf{T}}v$
slijede implikacije:
\begin{itemize}[nosep,leftmargin=*]
  \item ako je $\mu_j^{\star}$ \textbf{velik}: $\popt$ znatno raste ako
        ,,stisnemo`` ograničenje $g_j(x) \leq 0$ na $g_j(x) \leq v_j$ gdje je $v_j < 0$;
  \item ako je $\mu_j^{\star}$ \textbf{malen}: $\popt$ se ne smanjuje znatno ako
        ,,relaksiramo`` ograničenje na $g_j(x) \leq v_j$ gdje je $v_j > 0$.
\end{itemize}
(Slajd navodi i analogne implikacije za množitelje $\lambda_i$ ograničenja
jednakosti, kojih u našem zadatku nema.)

\textbf{Zašto je ovo \emph{globalna} tvrdnja:} ovo je nejednakost, ne
aproksimacija, i vrijedi za \emph{svaku} veličinu perturbacije --- ne samo za
malu. Cijena je što daje samo jednostranu granicu.
\end{izvor}

\begin{formula}{Globalna donja granica}
\begin{equation}
\popt(v) \;\geq\; \popt(0) - \muopt{}\trans v
\qquad\text{za svaki } v .
\label{eq:z1f-globalna}
\end{equation}
Uz $\muopt = (404/33,\; 0)$ i $\popt(0) = 2000/33$:
\begin{equation}
\popt(v) \;\geq\; \frac{2000}{33} - \frac{404}{33}\,v_1 .
\end{equation}
\end{formula}

\section{Lokalna osjetljivost}
\label{sec:z1f-lokalna}

\begin{izvor}{Predavanje 4 --- Lagrangeova dualnost, slajd 26/26 (PDF str.~34)}
\slajd[0.86]{p04-lokalna-osjetljivost}

\textbf{Transkripcija.} \emph{Lokalna osjetljivost.} Pretpostavimo jaku dualnost
za neperturbiran problem i neka su $\lambda^{\star}$, $\muopt$ optimalne dualne
varijable primarnog problema. Ako je $\popt(u,v)$ diferencijabilno u $(0,0)$,
tada vrijedi
\[
\lambda_i^{\star} = -\frac{\partial \popt(0,0)}{\partial u_i},
\qquad
\mu_i^{\star} = -\frac{\partial \popt(0,0)}{\partial v_i}.
\]

\textbf{Dvije pretpostavke koje se ne smiju prešutjeti:} jaka dualnost (kod nas
osigurana Slaterovim uvjetom) i \emph{diferencijabilnost} $\popt$ u nuli. Druga
zna pasti --- $\popt(v)$ je konveksna po $v$, ali može imati lom ondje gdje se
mijenja aktivni skup.
\end{izvor}

\begin{formula}{Osjetljivost --- ,,cijena u sjeni``}
\begin{equation}
\left.\frac{\partial \popt}{\partial v_i}\right|_{v = 0} = -\mu_i^{\star} .
\label{eq:z1f-osjetljivost}
\end{equation}
Lagrangeov množitelj je, do na predznak, \textbf{cijena} $i$-tog ograničenja:
koliko se optimalna vrijednost popravi po jedinici popuštanja.
\end{formula}

\begin{korak}[Brojevi za naš zadatak]
Uz $\muopt = (404/33,\; 0)$:
\begin{equation}
\frac{\partial \popt}{\partial v_1} = -\frac{404}{33} = -12{,}242424,
\qquad
\frac{\partial \popt}{\partial v_2} = 0 .
\end{equation}
Čitanje:
\begin{itemize}[leftmargin=*]
  \item \textbf{Prvo ograničenje.} Popuštanje za $\Delta b_1 > 0$ smanjuje
        $\popt$ za približno $12{,}2424 \cdot \Delta b_1$. Stezanje ga za isto
        toliko povećava. Ovo je ograničenje koje ,,stoji`` --- ono drži rješenje.
  \item \textbf{Drugo ograničenje.} U prvom redu nema nikakvog utjecaja, jer je
        $g_2$ u optimumu neaktivno pa je $\mu_2^{\star} = 0$. Optimum ga
        jednostavno ne dodiruje.
\end{itemize}
\end{korak}

\subsection{Dokle vrijedi ,,nema utjecaja``}

\begin{korak}[Rezerva do promjene aktivnog skupa]
Tvrdnja $\partial\popt/\partial v_2 = 0$ vrijedi samo dok $g_2$ ostaje neaktivno.
Rezerva iznosi
\begin{equation}
\abs{g_2(\xopt)} = \frac{49}{33} = 1{,}484848 .
\end{equation}
Ako se drugo ograničenje stegne za više od toga --- dakle $v_2 < -49/33$ ---
ono postaje aktivno, aktivni skup se mijenja, $\mu_2$ postaje pozitivan, i od tog
trenutka perturbacija \emph{jest} skupa. Baš na toj granici $\popt(v)$ ima lom i
nije diferencijabilna, pa pretpostavka s Predavanja 4 prestaje vrijediti.
\end{korak}

\begin{zamka}[Zamka: ,,$\mu_2 = 0$ znači da ograničenje ne treba``]
Ne znači. Ograničenje $g_2$ i dalje je dio problema i i dalje omeđuje dopustivi
skup. Znači samo da u \emph{ovom} optimumu nije aktivno, pa infinitezimalna
promjena njegove desne strane ne mijenja rješenje. Za promjenu veću od rezerve
$1{,}4848$ ograničenje postaje ono koje veže. Reci to eksplicitno --- to je
razlika između odgovora i naučenog odgovora.
\end{zamka}

\begin{zamka}[Zamka: predznak]
Formula ima minus: $\partial\popt/\partial v_i = -\mu_i^{\star}$, a $\mu_i^{\star} \geq 0$,
pa je derivacija $\leq 0$. Smisao: popuštanje ograničenja ($v_i > 0$) povećava
dopustivi skup, a minimum po većem skupu ne može biti veći. Ako ti u odgovoru
ispadne da popuštanje \emph{povećava} $\popt$, negdje si izgubio predznak.
\end{zamka}

\section{Numerička provjera}
\label{sec:z1f-provjera}

\begin{seminar}{Seminar v1.2, §2.7 --- ,,(f) Osjetljivost optimalne vrijednosti``}
Perturbacija desnih strana ograničenja uvodi se vektorom $v = [v_1, v_2]^{\mathsf{T}}$:
\[
p^*(v) = \min \left\{ \tfrac{1}{2} x^{\mathsf{T}} Q x + c^{\mathsf{T}} x \ \middle| \ Gx - h \leq v \right\},
\]
pri čemu $v_i$ odgovara pomaku desne strane $b_i$ $i$-tog ograničenja. Za
konveksan problem sa zadovoljenim Slaterovim uvjetom vrijedi
\[
\frac{\partial p^*}{\partial v_i}\bigg|_{v = 0} = -\mu_i^*.
\]
Uz $\mu^* = [404/33,\ 0]^{\mathsf{T}}$ slijedi
\[
\frac{\partial p^*}{\partial v_1} = -\frac{404}{33} = -12.242424, \qquad
\frac{\partial p^*}{\partial v_2} = 0.
\]
Optimalna vrijednost osjetljiva je samo na perturbaciju prvog ograničenja:
opuštanje ograničenja $g_1$ za $\Delta b_1 > 0$ smanjuje $p^*$ za približno
$12.2424 \, \Delta b_1$. Perturbacija drugog ograničenja nema utjecaja u prvom
redu jer je $g_2$ u optimumu neaktivno, tj.\ $\mu_2^* = 0$; to vrijedi sve dok je
perturbacija dovoljno mala da se aktivni skup ne promijeni (rezerva do aktivacije
iznosi $|g_2(x^*)| = 1.484848$).

\begin{lstlisting}[style=matlabstyle]
dlt = 1e-4;
[~, f_v1] = quadprog(Q, c, G, h + [dlt; 0], [], [], [], [], [], opts);
[~, f_v2] = quadprog(Q, c, G, h + [0; dlt], [], [], [], [], [], opts);
dp1 = (f_v1 - f_qp)/dlt;
dp2 = (f_v2 - f_qp)/dlt;
\end{lstlisting}

\begin{lstlisting}[style=outputstyle]
dp*/dv1 numericki = -12.242364,  -mu1* = -12.242424,  odstupanje = 5.99e-05
dp*/dv2 numericki = 0.000000,  -mu2* = 0.000000,  odstupanje = 0.00e+00
\end{lstlisting}

Jednostrana numerička derivacija odstupa od $-\mu_1^*$ za $5.99 \cdot 10^{-5}$,
što je reda veličine samog koraka $\delta$ i u skladu s teorijskim rezultatom.
\end{seminar}

\begin{korak}[Zašto odstupanje nije nula]
Korištena je \emph{jednostrana} razlika, $\bigl(\popt(\delta) - \popt(0)\bigr)/\delta$,
čija je pogreška reda $O(\delta)$. Uz $\delta = 10^{-4}$ očekuje se odstupanje
reda $10^{-4}$ do $10^{-5}$ --- dobiveno je $5{,}99\cdot 10^{-5}$, dakle točno
onoliko koliko metoda daje. Da je korištena centralna razlika,
$\bigl(\popt(\delta) - \popt(-\delta)\bigr)/(2\delta)$, pogreška bi bila $O(\delta^2)$.
Odstupanje, dakle, nije greška u teoriji nego svojstvo odabrane numeričke sheme.
\end{korak}

\section{Isti podzadatak na vježbama}
\label{sec:z1f-vjezbe}

\begin{izvor}{Vježbe 3 --- Primjer 2, slajd 34/36 (PDF str.~41)}
\slajd[0.86]{v3-primjer2-osjetljivost}

\textbf{Transkripcija.} \emph{e) Što možete reći o osjetljivosti optimalne
vrijednosti na perturbacije u ograničenjima?} Budući da je
$\mu^{\star} = [0,\ 2,\ 0,\ 2]\trans$, imamo
\[
\frac{\partial \popt(0,0,0,0)}{\partial v_1} = 0,\quad
\frac{\partial \popt}{\partial v_2} = -2,\quad
\frac{\partial \popt}{\partial v_3} = 0,\quad
\frac{\partial \popt}{\partial v_4} = -2 .
\]

\textbf{Struktura odgovora je identična našoj:} pročitaj množitelje, stavi minus,
i to je cijela osjetljivost. Ograničenja s nultim množiteljem su neaktivna i ne
utječu. Razlika je samo u tome što u profesorovu primjeru \emph{dva} ograničenja
imaju množitelj različit od nule, a kod nas jedno.
\end{izvor}

\section{Modelski odgovori}
\label{sec:z1f-odgovori}

\begin{pitanje}{Što biste rekli inženjeru koji pita treba li uložiti u popuštanje ograničenja?}
\textbf{Odgovor.} Da se isplati ulagati isključivo u prvo ograničenje, i to sve
dok je cijena jedinice popuštanja manja od $12{,}24$ u jedinicama funkcije cilja.
Ulaganje u drugo ograničenje ne donosi ništa --- ono nije aktivno, pa se rješenje
ne bi pomaknulo. Množitelj je doslovno \emph{cijena u sjeni}: koliko bi platio za
jedinicu labavosti.
\end{pitanje}

\begin{pitanje}{Je li $\popt(v)$ linearna funkcija od $v$?}
\textbf{Odgovor.} Nije, nego konveksna i po dijelovima glatka. Derivacija
$-\muopt$ vrijedi lokalno, dok se aktivni skup ne promijeni. Nejednakost
$\popt(v) \geq \popt(0) - \muopt{}\trans v$ vrijedi globalno, i geometrijski
znači da je afina funkcija na desnoj strani \emph{potporna ravnina} konveksne
funkcije $\popt$ u točki $v = 0$.
\end{pitanje}

\begin{pitanje}{Kako biste odredili osjetljivost da je problem nekonveksan?}
\textbf{Odgovor.} Ovako ne bih. Izvod se oslanja na jaku dualnost, koja kod
nekonveksnih problema općenito ne vrijedi --- ostaje samo slaba dualnost, pa se
iz množitelja dobiva jednostrana granica, ne derivacija. Osjetljivost bi se onda
procjenjivala numerički, ponovnim rješavanjem perturbiranog problema, uz svijest
da rješavač može završiti u drugom lokalnom minimumu.
\end{pitanje}

\begin{izgovor}
,,Prvo bih razjasnio oznaku: $b_1$ i $b_2$ iz teksta zadatka su desne strane
ograničenja, minus deset i tri. Perturbaciju uvodim vektorom ve, tako da
ograničenja postanu ge iks minus ha manje ili jednako ve.

Budući da je problem konveksan i Slaterov uvjet zadovoljen, vrijedi jaka
dualnost, pa parcijalna derivacija optimalne vrijednosti po $v_i$ u nuli iznosi
minus mi i-to optimalno. Kod nas je mi jednak četiristo četiri kroz trideset tri i
nula, pa je osjetljivost po prvom ograničenju minus dvanaest zarez dva četiri, a
po drugom nula.

Inženjerski: samo prvo ograničenje ima cijenu. Popuštanje prvog ograničenja za
jedinicu smanjuje optimalnu vrijednost za dvanaest zarez dva četiri. Drugo
ograničenje je neaktivno, ima rezervu jedan zarez četiri osam, i sve dok
perturbacija ne premaši tu rezervu, na optimum ne utječe. Preko te granice
aktivni skup se mijenja i optimalna vrijednost ima lom.

Provjerio sam i numerički, konačnim razlikama s korakom deset na minus četiri:
dobio sam minus dvanaest zarez dva četiri dva tri šest četiri, uz odstupanje šest
puta deset na minus pet --- što je red veličine samog koraka, jer je razlika
jednostrana.``
\end{izgovor}
